% Sections 9 and 10 of the white paper, LaTeX module.
% To be \input into the master document. Requires amsmath.
% Plain-language commentary is set with \plain{...}; define e.g.
% \newcommand{\plain}[1]{\begin{quote}\itshape #1\end{quote}}

\section{Context and the cosmological beta}
\label{sec:context}

\subsection{Prior art and positioning}

The algebraic backbone of this section belongs to an established corpus. The mass-independence of the all-horizon entropy product was proven family by family by Cveti\v{c}, Gibbons, and Pope \cite{CGP2011}, and the all-horizon entropy sum, with virtual horizons explicitly included, was shown by Wang, Xu, and Meng \cite{WXM2014} and by Du and Tian \cite{DuTian} to depend only on the cosmological constant and matter couplings. Xu, Wang, and Meng \cite{XWM2014} derived first-law and Smarr structure on the virtual horizons themselves. Visser \cite{Visser2013} exhibited the Schwarzschild--de Sitter counterexample for the product. Most importantly, Zhang and Gao \cite{ZhangGao} reduced the whole question to two criteria that follow from the per-horizon first law, $\partial S_i/\partial M = 1/T_i$: the entropy product over a set of horizons is mass-independent if and only if $\sum_i 1/(T_i S_i)=0$ over that set, and the sum if and only if $\sum_i 1/T_i = 0$, together with a Vandermonde lemma and a sharp Laurent-exponent criterion on $f(r)$ that decides both without solving for the roots.

This section adds no new identity to that corpus. What it adds is a change of question. The literature asks which relations are mass-independent over which sets of roots. We ask what the gap between the physical subset and the complete root set does to an observer who only sees the physical subset, and we show that the gap behaves exactly like an omitted variable in a regression: what looks like broken universality, or like noise, is deterministic context that has been left out of the accounting.

\plain{Previous work established which combinations of horizons give clean, mass-free numbers, provided you count the ghost horizons too. Our question is different: what does the world look like if you refuse to count the ghosts? The answer is that it looks noisy and broken, and that all of the noise is fake.}

\subsection{The decomposition: local invariant equals universal over context}

Reissner--Nordstr\"om--de Sitter has the horizon quartic
\begin{equation}
r^4 - \ell^2 r^2 + 2M\ell^2 r - Q^2\ell^2 = 0,
\label{eq:rnds-quartic}
\end{equation}
with three physical roots $r_-, r_+, r_c$ and one negative virtual root $r_o$. Vieta gives the complete-set product as the constant term, so the complete-set entropy product is mass-free, consistent with the corpus above:
\begin{equation}
S_{\mathrm{all}} \equiv \pi^4 (r_- r_+ r_c\, r_o)^2 = \pi^4 Q^4 \ell^4.
\label{eq:Sall}
\end{equation}
The physical-subset product is then forced into the form
\begin{equation}
S_{\mathrm{phys}} \equiv \pi^3 (r_- r_+ r_c)^2 = \frac{\pi^3 Q^4 \ell^4}{\Sigma^2},
\qquad \Sigma \equiv r_- + r_+ + r_c = -r_o,
\label{eq:Sphys}
\end{equation}
where the second equality in the definition of $\Sigma$ is Vieta again (the quartic has no cubic term). Equation~\eqref{eq:Sphys} is the organizing statement of the section: the local invariant equals the universal invariant divided by a context factor, and the context factor $\Sigma$, the sum of the physical radii, is identically the virtual root. All the mass dependence of the physical product is carried by $\Sigma$. Verified exactly: at $Q=1$, $\ell=10$, $S_{\mathrm{all}}/\pi^4 = 10000$ for every $M$, while $S_{\mathrm{phys}}/\pi^3$ runs through $80.55, 79.51, 78.51, 77.55, 76.61$ as $M$ runs from $1.3$ to $1.7$, with the identity $Q^4\ell^4/\Sigma^2 = (r_-r_+r_c)^2$ holding to machine precision.

\plain{The clean number is always there. When you only count the horizons you can see, you are looking at the clean number divided by the size of the box, and the box grows with the mass. The mess is the box, not the law.}

\subsection{The beta and the regression}

Because \eqref{eq:Sphys} is a power law, the language of factor regression applies literally. Write $\log S_{\mathrm{phys}} = 4\log Q + 4\log\ell - 2\log\Sigma + \mathrm{const}$: the exponents are betas, and $\Sigma$ is the context factor whose omission biases everything else. Across a cloud of 400 black holes with independently drawn $(M,Q,\ell)$, regressing $\log S_{\mathrm{phys}}$ on the full context $(\log Q, \log\ell, \log\Sigma)$ recovers the betas $(4,4,-2)$ exactly, with $R^2 = 1$ to machine precision. Regressing on $\log\Sigma$ alone yields $R^2 = 0.12$. The residual scatter of the naive regression is not statistical noise: it is the signature of omitted variables, and it vanishes identically when the accounting is completed. The same lesson holds for the breakings established earlier in the paper: the Kerr--AdS breaking of the diagonal invariant is an exact analytic function of the pressure parameter $\varepsilon$ (a degree-five polynomial fit reaches $R^2 = 1.0000000000$), and the scattered Myers--Perry $d=6$ values collapse onto a single deterministic surface $F(J_a/M^2, J_b/M^2)$ once scaling is quotiented out. Nothing in this landscape is random.

\plain{In finance, if a stock looks wildly volatile, the first suspicion is not that the world is random but that your model is missing a factor. Here we can prove it: add the missing factor, the ghost horizon, and the volatility drops to exactly zero. Not approximately. Exactly.}

\subsection{Completability and its boundary}

The Vieta carrier identity, physical product times virtual product equals a mass-free constant, was checked explicitly for $d=4$ through $7$ in the charged de Sitter family. The $d=5$ case is special in a way that must be distinguished from the known dimension statements of the corpus (which concern the mass-independence of complete-set sums for $d>4$): in $d=5$ the reduced polynomial in $x=r^2$ has degree equal to the number of physical horizons, so the physical product is already mass-free by itself, with no virtual completion needed. We call this self-completion, and $d=5$ is the only dimension in the range examined where it occurs.

Rotation marks the boundary of the mechanism. For Kerr--de Sitter, the entropy factor is $r_i^2 + a^2$, and the complete-set product evaluates by Vieta to
\begin{equation}
\prod_{i=1}^{4}\left(r_i^2 + a^2\right) = P(ia)\,P(-ia) = 4M^2 a^2 \ell^4,
\label{eq:rotation-boundary}
\end{equation}
mass-dependent even over all four roots. Charge is context-completable; rotation is not. The frame-dragging factor leaves an irreducible beta that no accounting of virtual horizons removes. This is the same charge--rotation asymmetry that reappears, with its exact mechanism, in the branch-split anomaly of Section~\ref{sec:virtual} (equation~\eqref{eq:Wsplit}): there the quantity $\Sigma_W$ detects the scale anomaly that background charge induces and rotation never does, and here rotation is the sector whose context cannot be completed at all. The two statements are faces of one asymmetry and their common origin is open.

\subsection{The junction with the wave equation}

The completion statement of this section has a wave-equation face, developed in Section~\ref{ssec:monodromy}: the monodromy exponent of a mode $(\omega, m)$ at horizon $i$ is the first-law entropy shift per absorbed quantum, so the total exponent over the complete root set is the $(\omega\,\partial_M + m\,\partial_J)$ derivative of the entropy sum, and its vanishing is exactly the sum universality of Wang, Xu, and Meng, by the criterion of Zhang and Gao. The two literatures state the same closure without citing each other; the virtual horizons are the agents of it in both. The identities are theirs; the junction, and the identification of the virtual roots as the unfolded irregular singularity carrying quantized curvature charges, are the contribution of this paper.

\plain{The bookkeeping that makes the entropy numbers clean and the bookkeeping that makes waves scatter consistently around the black hole turn out to be the same bookkeeping, kept by the same ghosts. Two communities discovered it separately, in different words, and never noticed.}

% ---------------------------------------------------------------------------

\section{Virtual sheets and quantized curvature charges}
\label{sec:virtual}

\subsection{The complete cover and a failed conjecture}

Turning on the cosmological constant promotes the horizon condition from a quadratic to a quartic. For Reissner--Nordstr\"om--AdS at fixed AdS radius $\ell$,
\begin{equation}
\frac{r^4}{\ell^2} + r^2 - 2Mr + Q^2 = 0,
\label{eq:rnads-quartic}
\end{equation}
so the two-sheeted physical cover of the earlier sections is embedded in a four-sheeted branched cover of parameter space. Two roots are the physical horizons $r_\pm$; the other two form a complex conjugate pair, the virtual sheets. Section~\ref{sec:context} established that Vieta symmetric functions are context-complete on the full root set. The four-face equivalence ties the diagonal invariant $R_+S_+ + R_-S_-$ to product mass-independence. Assembling the two suggests a natural conjecture: the diagonal invariant, broken at finite pressure on the physical pair, should be restored as a Galois trace over the complete cover,
\begin{equation}
\sum_{i=1}^{4} R_i S_i \stackrel{?}{=} \mathrm{const},
\label{eq:galois-trace}
\end{equation}
where each branch entropy $S_i(M,Q) = \pi r_i^2$ defines its own Ruppeiner geometry (metric $g = -\mathrm{Hess}\,S$, the convention calibrated so that Kerr yields exactly $-1$) and the conjugate pair contributes a real sum.

This conjecture is false. The four-root sum is real, as it must be, but not constant: at $(M,Q,\ell) = (1, 0.5, 10)$ it equals $-3.211$, at $(1, 0.5, 5)$ it equals $-2.279$, at $(2, 0.7, 8)$ it equals $-0.729$. Galois invariance guarantees only that the trace is a rational function of the coefficients, not that it is constant. We record the refutation and turn to what the decomposition revealed instead.

\plain{We hoped that adding up the invariant over all four sheets, real and virtual, would give back a conserved number. It does not. But splitting the sum into its physical and virtual halves uncovered something better.}

\subsection{The virtual charges}

Decompose the trace into the physical pair and the virtual pair and take the zero-pressure limit $\ell \to \infty$. The physical pair returns to its flat value, $0$ for the $(M,Q)$ family and $-1$ for the $(M,J)$ family, which independently validates the numerical pipeline. The virtual pair converges to an exact rational number that depends on nothing except which two-plane of thermodynamic fluctuations defines the Ruppeiner geometry.

For electric fluctuations, coordinates $(M,Q)$, including on a Kerr--Newman--AdS background with arbitrary fixed rotation,
\begin{equation}
\lim_{\ell\to\infty}\,\bigl(R_v S_v + R_{\bar v} S_{\bar v}\bigr)\Big|_{(M,Q)} = -\frac{27}{8},
\label{eq:charge-electric}
\end{equation}
independent of the background values of $M$, $Q$, and $J$. For angular fluctuations, coordinates $(M,J)$, on a background of fixed charge $Q$,
\begin{equation}
\lim_{\ell\to\infty}\,\bigl(R_v S_v + R_{\bar v} S_{\bar v}\bigr)\Big|_{(M,J)}
= -\frac{25}{8} + \frac{5}{16}q^2 + \frac{1}{16}q^4
= \frac{(q^2+10)(q^2-5)}{16},
\qquad q \equiv \frac{Q}{M},
\label{eq:charge-angular}
\end{equation}
independent of $J$ and of overall scale, deformed only by the background charge. Both statements are exact symbolic theorems (Sec.~\ref{ssec:closedforms}); they were first found numerically, with Richardson extrapolation in $1/\ell^2$ agreeing with the rational values to ten or more digits and a predictive out-of-grid test at $q=0.7$ confirmed to eleven digits. Equation~\eqref{eq:charge-electric} holds sheet by sheet: each virtual sheet carries exactly $-27/16$, a real number, with every negative power of $1/\ell$ vanishing identically.

\plain{The two ghost horizons, the complex roots nobody assigns a temperature to, carry a fixed curvature charge. In the electric frame that charge is minus twenty-seven eighths, always. Rotate the background as much as you like, it does not move. In the angular frame the charge starts at minus twenty-five eighths and is pushed by background charge along an exact polynomial, and by nothing else.}

\subsection{The rigidity hierarchy}

The virtual charges are more rigid than the physical integers. In the $(M,Q)$ frame at fixed $J\neq 0$, the flat-limit physical pair sum is not even constant: it takes the values $0.404$, $0.344$, $0.504$ at $(M,Q) = (1,0.4), (1.5,0.6), (1,0.7)$ with $J=0.3$. The virtual pair at those same points is $-27/8$ to full precision. The completion sector holds its value while the physical sector drifts with the background. This is the curvature-level image of the context decomposition of Section~\ref{sec:context}: $R\cdot S$ on the virtual sheets is a pure number of the fluctuation frame, a zero-beta quantity, while all context sensitivity is carried by the physical sheets. The virtual completion is the zero-beta sector of the branched cover.

The charge--rotation asymmetry of Section~\ref{sec:context} also reappears, with an inversion we flag honestly: background rotation is inert everywhere (proven exactly in the electric frame), while background charge is the sole agent of deformation, acting only on the angular frame through the even polynomial of \eqref{eq:charge-angular}. We note that $Q$ enters the deformation starting at $q^2$, whereas the level mismatch of the four-face analysis carries charge as $Q^4/4$ inside $N_L - N_R$; the two exponents do not match naively and any CFT reading of the virtual charges must account for this.

\subsection{Analytic derivation}
\label{ssec:analytic}

The flat limit is controlled by the asymptotics of the virtual root. Writing $x = 1/\ell$, the virtual branch of \eqref{eq:rnads-quartic} is
\begin{equation}
r_v = \frac{i}{x} - M + \frac{i(3M^2 - Q^2)}{2}\,x + 2M(2M^2 - Q^2)\,x^2 + O(x^3),
\label{eq:virtual-series}
\end{equation}
with the conjugate sheet obtained by $i \to -i$. The branch entropy $S_v = \pi r_v^2$ diverges as $-\pi\ell^2$, while the leading Hessian metric is the constant Lorentzian form $\mathrm{diag}(4\pi, -2\pi)$; curvature therefore enters at order $\ell^{-2}$ and the product $R\cdot S$ is finite in the limit. The exact evaluation proceeds by implicit differentiation of \eqref{eq:rnads-quartic} through third order (the Brioschi formula for a Hessian metric requires nothing beyond third order, since its second-derivative block cancels identically for any Hessian metric), substitution of the Laurent series \eqref{eq:virtual-series} extended by Newton iteration in truncated Laurent-series arithmetic over the Gaussian rationals, and extraction of the constant term of $R\cdot S$, all negative powers vanishing identically.

\subsection{Closed-form theorems}
\label{ssec:closedforms}

Carrying the machinery of Sec.~\ref{ssec:analytic} with the parameters left symbolic (mass set to one without loss of generality by scale invariance) gives, per virtual sheet: in the electric frame with generic symbolic $Q$,
\begin{equation}
R_v S_v = -\frac{27}{16};
\label{eq:persheet-electric}
\end{equation}
in the angular frame with generic symbolic $Q$,
\begin{equation}
R_v S_v = -\frac{25}{16} + \frac{5}{32}q^2 + \frac{1}{32}q^4,
\label{eq:persheet-angular}
\end{equation}
a real number, so each sheet of the conjugate pair carries exactly half the pair charge.

\subsection{The branch-level level map and the general split theorem}
\label{ssec:numap}

For Kerr, the individual products admit an exact closed form. With $\nu \equiv \sqrt{N_R/N_L} = \sqrt{1 - J^2/M^4}$, which is also the entropy asymmetry $(S_+-S_-)/(S_++S_-)$,
\begin{equation}
R_\pm S_\pm = -\frac{1}{2} \pm \frac{N_R - 2N_L}{2\sqrt{N_L N_R}}
= -\frac{1}{2} \pm \left(\frac{\nu}{2} - \frac{1}{\nu}\right).
\label{eq:numap}
\end{equation}
The sum is identically $-1$; the deck involution $\nu \to -\nu$ exchanges the two expressions. Equation~\eqref{eq:numap} is the precise map from the curvature invariant to the Kerr/CFT level structure at the level of individual branches: each horizon carries a universal charge of $-1/2$ plus a level-mismatch term of opposite sign, and the topological $-1$ is the trace over the two-sheet cover. It matches the unified curvature of \.{A}man, Bengtsson, and Pidokrajt (their Eqs.~(35) and~(37)) on both branches exactly, an independent convention check.

The scope is settled by a general theorem. For any family of the deck-symmetric form $S_\pm = 2\pi\bigl(W(M) \pm \sqrt{W(M)^2 - J^2}\bigr)$ with $W$ an arbitrary function of the mass,
\begin{equation}
R_+S_+ + R_-S_- = -1 \quad\text{identically},
\qquad
R_\pm S_\pm = -\frac{1}{2} \pm D,
\quad
D = \frac{(3 - \Sigma_W)\,\nu}{2} - \frac{1}{\nu},
\quad
\Sigma_W \equiv \frac{W' W'''}{W''^2} + \frac{W'^2}{W\,W''}.
\label{eq:Wsplit}
\end{equation}
The integer is a property of the $\pm$ structure alone (the product $S_+S_- = 4\pi^2 J^2$ is automatically mass-free family-wide, so this is the four-face equivalence holding for the whole class). The scale-invariant quantity $\Sigma_W$ equals $2$ for every power law $W = m^p$, independently of $p$: the Kerr split is shared by all scale-covariant families, which is why equal-spin Myers--Perry in five dimensions ($p = 3/2$) follows the Kerr map exactly (verified at machine precision, alongside a forty-digit cross-validation of the outer branch against the closed form of Bravetti, Nettel, and Quevedo). Kerr--Sen has $W = M^2 - Q^2/2$, a shifted power law, giving $\Sigma_W = 2/(1 - q^2/2)$; the resulting deviation matches the measured values exactly. The trace is topological and blind to the level coordinate; the branch split measures the scale anomaly of the level coordinate, which background charge induces and rotation does not.

\subsection{Frame sensitivity of the geometric face}

For Kerr--Newman in the two-dimensional angular frame $(M,J)$ at fixed $Q$, the diagonal sum drifts continuously with $Q$ and approaches $-1$ only as $Q \to 0$, consistent with the charge-deformation coefficient established earlier in the paper, while the product face $S_+S_- = 4\pi^2(J^2 + Q^4/4)$ remains exactly mass-free. The geometric face of the equivalence is a property of a fluctuation frame; the product and CFT faces are frame-blind. The natural universality test is Kerr--Sen, where \eqref{eq:Wsplit} now supplies the exact split for all charge.

\subsection{Quantization pattern and a conjecture}

Per virtual sheet the two frame charges are $-27/16 = -3^3/16$ and $-25/16 = -5^2/16$; all exact values obtained, including the deformation coefficients $5/32$ and $1/32$ per sheet, are quantized in thirty-seconds per sheet. We flag the $3^3$ and $5^2$ pattern as an observation, not a claim. The two undeformed charges fit a one-parameter formula in the scaling weight $w$ of the second fluctuation coordinate ($w=1$ for $Q$, $w=2$ for $J$),
\begin{equation}
\text{pair charge} = -\frac{29}{8} + \frac{w}{4},
\label{eq:weight-conjecture}
\end{equation}
a conjecture determined by two points and falsifiable on Kerr--Sen (prediction $-27/8$ in its electric frame) and on five-dimensional Myers--Perry at fractional weight.

\subsection{Monodromy closure and the identity of the virtual sheets}
\label{ssec:monodromy}

Two exact statements tie the cover to the wave equation. For a neutral massless scalar on RN--AdS the monodromy exponent at each root is $\alpha_i = \omega r_i^2/\Delta'(r_i)$ up to normalization, and for Kerr--AdS it is $\alpha_i = [\omega(r_i^2 + a^2) - m a(1 + r_i^2/\ell^2)]/\Delta'(r_i)$. In both cases
\begin{equation}
\sum_{i=1}^{4} \alpha_i = 0 \quad\text{identically, for all } \omega \text{ and } m,
\label{eq:monodromy-closure}
\end{equation}
verified to forty digits and provable in one line by the residue theorem, since the AdS quartic leaves no residue at infinity. The physical content is sharper than folklore: the exponent at horizon $i$ is the first-law entropy shift per absorbed quantum, $\alpha_i = (\omega\,\partial_M + m\,\partial_J)S_i/4\pi$, so the total is the same derivative of the entropy sum. The closure \eqref{eq:monodromy-closure} is therefore exactly equivalent, via the criterion of Zhang and Gao, to the mass- and angular-momentum-independence of the all-horizon entropy sum of Wang, Xu, and Meng and Du--Tian: two literatures, black hole monodromy and entropy relations, state the same fact without citing each other, and the virtual horizons are the agents of the closure in both.

The flat-limit bookkeeping identifies the virtual sheets. As $\ell \to \infty$ the physical pair sum tends to $2M\omega$, the known flat value, while the virtual pair sum tends to exactly $-2M\omega$: precisely the monodromy content that flat space carries at its irregular singular point at infinity, as Stokes data. In AdS there is no irregular point; the two virtual roots sit at $r \approx \pm i\ell - M$ as regular singular points, and in the flat limit they run to infinity and merge, the standard confluence that manufactures an irregular singularity from two regular ones. The virtual sheets are the AdS unfolding of the flat-space irregular singularity, and the quantized charges \eqref{eq:persheet-electric} and \eqref{eq:persheet-angular} are thermodynamic-geometry data attached to the unfolded object.

\plain{In flat space the black hole's wave equation hides part of its structure at infinity, folded up in a mathematical object called an irregular singularity. Turning on the cosmological constant unfolds it into two visible points, and those two points are exactly the ghost horizons whose curvature charges we computed.}

The isomonodromy program gives this a precise target. The massless scalar equation on RN--AdS is Fuchsian with five regular singular points; its isomonodromic deformations form a two-time Garnier system, and the thermodynamic fluctuation plane $(M,Q)$ at fixed $\ell$ has exactly that dimension. The framework for these deformations, up to five regular points with Painlev\'e~VI and Garnier tau functions, is established in the isomonodromy school for black hole radial equations; the conjecture proper to this paper is that the Ruppeiner branch geometry pulls back structures of the Garnier system, with the branch entropies related to the logarithm of the isomonodromic tau function localized at each singular point, and that $-27/16$ appears as tau-function data of the confluence limit.

\subsection{Status within the program}

The naive completion of the diagonal invariant fails: pressure does not conserve the four-sheet trace. What the completion does instead is split the invariant into rigid, quantized pieces attached to the Galois orbits of the horizon polynomial, with the virtual orbit carrying frame-dependent pure numbers and the physical orbit carrying all the context. The four-face equivalence gave the physical integer four readings; the virtual charge is a fifth object, attached to the completion rather than to either horizon, and it currently has no CFT face. Whether the charges admit a level-structure or tau-function reading, and whether the inert role of $J$ here is the mirror image of its obstruction role in Section~\ref{sec:context}, are the questions this section leaves open.

\plain{We did not find the conserved total we were looking for. We found something stranger: the imaginary part of the black hole's family tree carries fixed rational numbers, more stable than anything on the real side, and we can now compute them exactly. What they count, nobody knows yet.}
